
\chapter{Introdução}
\label{sec1}

Problemas de otimização multi-objetivo (\textit{Multi-Objective Problems}) popularmante conhecidos como MOPs abrangem uma ampla variedade de desafios do mundo real encontrados na ciência e engenharia. Estes incluem o problema da mochila multidimensional, problemas de alocação de recursos e controle de tráfego, entre outros~\cite{citacaon2}. Tais problemas caracterizam-se por envolver objetivos conflitantes, tornando sua otimização uma tarefa de alta complexidade.

Formalmente, MOPs podem ser definidos como o processo de otimizar simultaneamente duas ou mais funções objetivo conflitantes sujeitas a determinadas restrições, expressos matematicamente como:

\begin{align}
\label{eq:mop_definition}
\min \quad & \mathbf{F}(\mathbf{x}) = [f_1(\mathbf{x}), f_2(\mathbf{x}), \ldots, f_m(\mathbf{x})]^T \\
\text{sujeito a} \quad & \mathbf{x} \in \Omega \nonumber
\end{align}

\noindent onde $\Omega$ representa o espaço de soluções viáveis, $\mathbf{x}$ é uma solução descrita como um vetor de decisão, e $\mathbf{F}(\mathbf{x})$ é um vetor de $m$ funções objetivo $f_i(\mathbf{x})$ aplicadas sobre a solução $\mathbf{x}$.

Uma abordagem amplamente utilizada para superar as limitações dos métodos de otimização clássicos no tratamento de problemas de otimização multimodal e multi-objetivo emergiu com a introdução de técnicas de computação evolutiva (\textit{Evolutionary Computation})~\cite{citacaon4} as EC. Neste contexto, abordagens metaheurísticas como os Algoritmos Evolutivos Multi-objetivo (\textit{Multi-Objective Evolutionary Algorithms}) ou MOEAs foram propostas. Entre estes, destacam-se o NSGA-II~\cite{deb2002fast} e o SPEA2~\cite{zitzler2001spea2}, que ganharam considerável popularidade na literatura especializada.

Os MOPs podem apresentar diferentes configurações de problemas, com fronteiras de Pareto desafiadoras de alcançar e preencher, mesmo com algoritmos evolutivos tradicionais. À medida que a escala dos problemas aumenta com um número mais significativo de objetivos e variáveis de decisão o desempenho dos algoritmos se deteriora substancialmente~\cite{citacaon3}. Alcançar resultados de otimização satisfatórios para problemas caracterizados por descontinuidades de fronteira, espaços de busca multimodais e características diversas de fronteira requer um gerenciamento cuidadoso do \textit{trade-off} entre convergência e diversificação~\cite{huband}.

Neste contexto, Camilo-Junior e Yamanaka~\cite{camilo-junior2011} introduziram o método de Fertilização \textit{In Vitro} (\textit{In Vitro Fertilization}) ou IVF como um mecanismo projetado para melhorar a taxa de convergência e promover maior diversidade em algoritmos evolutivos. Este método tem a capacidade de gerar soluções promissoras mediante a manipulação das informações genéticas das soluções mais promissoras encontradas, utilizando diversas estratégias de busca local.

Para avaliar adequadamente a capacidade de exploração e intensificação de algoritmos evolutivos multi-objetivo, diferentes \textit{benchmarks} foram desenvolvidos, buscando explorar distintas configurações de problemas e geometrias da fronteira de Pareto. O \textit{benchmark} ZDT~\cite{zitzler2000comparison} contém seis problemas que apresentam exemplos de diferentes geometrias de fronteira de Pareto, oferecendo desafios significativos para algoritmos evolutivos tradicionais devido às suas fronteiras de Pareto convexas, côncavas, descontínuas e multimodais. O \textit{benchmark} DTLZ~\cite{deb2005scalable} foi desenvolvido para superar algumas limitações do ZDT, oferecendo problemas escaláveis com um número arbitrário de objetivos e permitindo melhor análise de desempenho em otimização de muitos objetivos (\textit{many-objective optimization}). Posteriormente, o \textit{benchmark} WFG~\cite{huband2006review} introduziu um conjunto mais abrangente de problemas apresentando fronteiras não separáveis, degeneradas e de alta dimensão, fornecendo um conjunto completo de ferramentas para testar algoritmos sob condições diversas. Mais recentemente, o \textit{benchmark} MaF~\cite{cheng2017benchmark} foi especificamente projetado para avaliar algoritmos em problemas de grande escala, combinando alta dimensionalidade no espaço de decisão com múltiplos objetivos, simulando melhor os desafios encontrados em aplicações do mundo real.

Entre os algoritmos evolutivos multi-objetivo destacados na literatura, o algoritmo SPEA2~\cite{zitzler2001spea2} se destaca como uma melhoria significativa sobre o Algoritmo Evolutivo de Pareto de Força original (\textit{Strength Pareto Evolutionary Algorithm}) SPEA. O SPEA2 combina uma estratégia aprimorada de alocação de \textit{fitness}, uma técnica sofisticada de estimativa de densidade e um método melhorado de truncamento de arquivo. Apesar de sua reconhecida eficácia, o SPEA2 tem sido relativamente pouco explorado em certos domínios de otimização complexos, apresentando oportunidades para hibridização a fim de abordar desafios específicos. 

As forças comprovadas do algoritmo incluem a capacidade de alcançar excelente diversificação de soluções na fronteira de Pareto, manter uma abordagem elitista eficaz e preservar soluções não dominadas de alta qualidade mesmo sob complexidade crescente do problema. Pesquisas demonstraram que o SPEA2 supera algoritmos como o NSGA-II à medida que o número de objetivos aumenta~\cite{reference1}, particularmente em termos de uniformidade de distribuição de soluções e precisão de convergência~\cite{reference2}. Este desempenho superior em cenários multi-objetivo é crucial para aplicações do mundo real onde múltiplos critérios conflitantes devem ser simultaneamente otimizados. Além disso, o SPEA2 foi aplicado com sucesso em campos diversos como otimização de \textit{design} de construção modular~\cite{Chen2021}, problemas de composição de serviços \textit{web} conscientes de QoS~\cite{Cremene2016} e otimização de agendamento de bombas em sistemas de distribuição de água~\cite{Lopez2005}, confirmando sua versatilidade e adaptabilidade.

No entanto, apesar dessas vantagens, o SPEA2 pode beneficiar-se da hibridização para melhorar suas capacidades de exploração e acelerar a convergência em espaços de busca altamente restringidos. Ao integrar técnicas complementares com a estrutura robusta do SPEA2, é possível abordar as limitações do algoritmo na eficiência de busca local enquanto se preservam suas capacidades de busca global. Esta abordagem de hibridização oferece uma estratégia promissora para enfrentar problemas de otimização caracterizados por espaços de decisão complexos, múltiplos objetivos conflitantes e paisagens de restrições desafiadoras.

Estudos recentes destacam a contribuição de algoritmos híbridos que intercalam buscas locais no fluxo de algoritmos evolutivos. Abordagens como IVF/NSGA-II~\cite{sampaio2017ivf} e IVF/GDE3~\cite{sampaio2019ivf} demonstraram as vantagens de incorporar o método IVF multi-objetivo em algoritmos evolutivos consolidados como o NSGA-II e o GDE3. Esta integração melhora características específicas dos algoritmos enquanto contribui para um equilíbrio entre convergência e diversificação. Em configurações específicas de problemas, o desempenho do SPEA2 foi igual ou superior aos algoritmos evolutivos mencionados acima~\cite{citacaon1}.

O método IVF, originalmente proposto por Camilo-Junior e Yamanaka~\cite{camilo-junior2011}, funciona como um procedimento elitista que manipula informações genéticas de soluções promissoras. Ele opera selecionando indivíduos de elite (gametas) da população atual, criando embriões através de operações de cruzamento predefinidas, e então avaliando e potencialmente incorporando essas novas soluções na população. Este processo permite a exploração de regiões promissoras do espaço de busca enquanto mantém a diversidade, abordando um desafio fundamental na otimização multi-objetivo.

Com base em um exame detalhado do algoritmo evolutivo multi-objetivo SPEA2, do método IVF e dos algoritmos híbridos IVF/NSGA-II e IVF/GDE3, este trabalho propõe, pela primeira vez, o acoplamento do algoritmo IVF com o SPEA2, compondo assim o algoritmo híbrido IVF/SPEA2. Esta combinação inovadora visa aproveitar as forças de ambas as abordagens: o desempenho superior do SPEA2 em manter soluções uniformemente distribuídas ao longo da fronteira de Pareto e a capacidade do método IVF de gerar indivíduos de alta qualidade através de manipulação genética direcionada.

Portanto, as principais contribuições desta dissertação podem ser resumidas como segue:

\begin{enumerate}
    \item \textbf{Desenvolvimento do algoritmo híbrido IVF/SPEA2:} Acoplamento inédito do método IVF com o SPEA2 canônico para contribuir na geração de indivíduos promissores dentro do fluxo algorítmico. Tal acoplamento garante a preservação de ambos os comportamentos de exploração (\textit{exploration}) e explotação (\textit{exploitation}) ao longo do procedimento de otimização;
    
    \item \textbf{Avaliação experimental abrangente:} Condução de uma avaliação empírica extensiva do IVF/SPEA2 através de quatro suítes de \textit{benchmark} (ZDT, DTLZ, WFG e MaF) comparando-o com algoritmos populares na literatura (SPEA2, NSGA-II, NSGA-III, MOEA/D) e hibridizações recentes do SPEA2 (MFOSPEA2 e SPEA2+SDE), demonstrando desempenho superior em cenários diversos de otimização multi-objetivo.
\end{enumerate}

O restante desta dissertação está estruturado para guiar o leitor através do trabalho desenvolvido da seguinte forma: o Capítulo~\ref{sec:Background} apresenta os conceitos essenciais dos algoritmos base e fundamentação teórica, seguido por uma revisão de trabalhos relacionados no Capítulo~\ref{sec:RW}. O núcleo da proposta, detalhando o acoplamento do método IVF ao algoritmo SPEA2, é apresentado no Capítulo~\ref{sec:Proposal}. Subsequentemente, o Capítulo~\ref{sec:Experiments} descreve a metodologia experimental e os testes conduzidos para avaliar a abordagem proposta. Uma análise aprofundada dos resultados, juntamente com discussões e padrões observados, é fornecida no Capítulo~\ref{sec:results}. Finalmente, o Capítulo~\ref{sec:discussion_conclusions} apresenta as conclusões e delineia direções potenciais para pesquisas futuras.